\chapter{Học Thêm Không Cứu Nổi Sự Phân Tán}
\markboth{Học Thêm Không Cứu Nổi Sự Phân Tán}{Học Thêm Không Cứu Nổi Sự Phân Tán}

\begin{center}
	\textit{\color{bookprimary}``Lịch học thêm kín không tự động biến một đứa trẻ thành người có nội lực.''}

	\textit{\color{bookprimary}``A packed tutoring schedule does not automatically turn a child into a person with inner strength.''}
\end{center}

\ngancach

\begin{tcolorbox}[colback=bookprimary!6,colframe=bookprimary!35,arc=5pt,title={\bfseries Góc Suy Nghĩ},fonttitle=\bfseries\color{bookprimary}]
Nhiều phụ huynh thấy con yếu là cho học thêm.
\textit{(Many parents see weakness in their child and respond by sending them to extra classes.)}

Nhiều học sinh thấy mình chênh vênh là đi học thêm như một cách chống hoảng.
\textit{(Many students feel unstable and attend tutoring as a way to quiet panic.)}

Học thêm không xấu.
\textit{(Tutoring is not inherently bad.)}

Nhưng khi đầu óc đã nát vì thiếu tập trung, ngủ ít, cầm điện thoại nhiều và học đối phó, thì thêm giờ chưa chắc thêm chất.
\textit{(But when the mind is already worn down by poor focus, little sleep, heavy phone use, and defensive studying, more hours do not necessarily mean more substance.)}
\end{tcolorbox}

\section{Lịch Kín Làm Người Lớn Yên Tâm}
\begin{truyen}{Lịch Kín Làm Người Lớn Yên Tâm}{A Packed Schedule Reassures Adults}
\chuhoa{C}{ó phụ huynh nhìn thời khóa biểu kín mít của con rồi thở phào.}
\textit{(Some parents look at their child's packed schedule and feel relieved.)}

Sáng học chính, chiều học thêm, tối làm bài, cuối tuần luyện đề.
\textit{(School in the morning, tutoring in the afternoon, homework at night, mock exams on weekends.)}

Nhìn rất giống chăm chỉ.
\textit{(It looks very much like diligence.)}

\teacherpair{Lịch kín đôi khi chỉ làm người lớn bớt lo, chứ chưa chắc làm học sinh bớt rỗng. Nếu con đi hết nơi này tới nơi khác mà đầu óc vẫn sống trong mù mờ, thì cái được lấp chỉ là lịch, không phải năng lực.}{``A packed schedule sometimes only makes adults less anxious; it does not necessarily make students less hollow. If a child moves from one class to another while the mind remains foggy, then what gets filled is only the calendar, not capability.''}
\end{truyen}

\section{Đi Học Thêm Nhưng Trong Đầu Vẫn Lướt}
\begin{truyen}{Đi Học Thêm Nhưng Trong Đầu Vẫn Lướt}{Attending Tutoring While the Mind Keeps Scrolling}
\chuhoa{C}{ó học sinh ngồi trong lớp thêm nhưng tay cứ ngứa điện thoại.}
\textit{(Some students sit in tutoring class while their hands keep itching for the phone.)}

Nếu không cầm máy thì đầu cũng nhảy như đang lướt dở.
\textit{(Even without touching the phone, the mind jumps around as if still mid-scroll.)}

\teacherpair{Học thêm không cứu nổi một đầu óc đã quen bị xé vụn. Muốn học được, trước hết phải sửa cái nền của sự chú ý. Nếu không, em sẽ chỉ mang một cơ thể tới lớp còn tâm trí vẫn ở trong dòng cuộn vô tận.}{``Tutoring cannot rescue a mind that has grown used to being torn into fragments. To learn again, the foundation of attention must be repaired first. Otherwise only your body arrives in class while your mind stays trapped in the endless feed.''}
\end{truyen}

\section{Thêm Thầy Không Thay Được Tự Học}
\begin{truyen}{Thêm Thầy Không Thay Được Tự Học}{More Teachers Cannot Replace Self-Study}
\chuhoa{N}{hiều bạn nghĩ}: cứ đi học thêm nhiều thì sẽ khá lên.
\textit{(Many students think that attending more tutoring classes will automatically make them better.)}

Không hẳn.
\textit{(Not necessarily.)}

\teacherpair{Thầy giỏi có thể kéo em đi một đoạn. Nhưng đoạn đường quyết định vẫn là lúc em ngồi một mình với vở, với lỗi sai, với sự nản và với ý chí của chính mình. Người nào không tập tự học thì sớm muộn cũng nghiện người dẫn.}{``A good teacher can pull you forward for a while. But the decisive stretch is when you sit alone with your notebook, your mistakes, your frustration, and your own will. Whoever never trains self-study eventually becomes dependent on being led.''}
\end{truyen}

\section{Học Để Đỡ Sợ, Không Phải Để Hiểu}
\begin{truyen}{Học Để Đỡ Sợ, Không Phải Để Hiểu}{Studying to Reduce Fear, Not to Understand}
\chuhoa{C}{ó những buổi học thêm mà mục tiêu thật không phải là hiểu bài.}
\textit{(There are tutoring sessions whose real goal is not understanding.)}

Mục tiêu là để bớt sợ rớt, bớt sợ thua, bớt sợ bị chê.
\textit{(The goal is to fear failure less, fear losing less, fear criticism less.)}

\teacherpair{Nỗi sợ có thể đẩy em tới lớp. Nhưng nếu cả quá trình học chỉ do sợ kéo đi, em rất khó có niềm vui học thật và cũng rất khó bền. Học vì sợ lâu ngày làm con người kiệt mà vẫn không sâu.}{``Fear can push you into the classroom. But if the whole learning process is driven only by fear, it becomes very hard to discover real joy in learning and very hard to endure. Studying from fear exhausts a person without making them deep.''}
\end{truyen}

\section{Lớp Thêm Thành Nơi Chuyền Mẹo}
\begin{truyen}{Lớp Thêm Thành Nơi Chuyền Mẹo}{Tutoring Class Becomes a Place for Tricks}
\chuhoa{C}{ó lớp thêm không nuôi tư duy mà chỉ chuyền mẹo làm nhanh.}
\textit{(Some tutoring classes do not cultivate thinking; they only pass around shortcuts.)}

Điểm có thể lên một lúc.
\textit{(Scores may rise for a while.)}

Nhưng nền rất mỏng.
\textit{(But the foundation remains thin.)}

\teacherpair{Mẹo không có tội nếu đi sau hiểu biết. Nhưng khi mẹo thế chỗ tư duy, học sinh bắt đầu quen sống bằng lối tắt. Lâu dần, các em không còn muốn hiểu nữa, chỉ muốn vượt qua nhanh. Đó là kiểu thành tích nghèo nàn nhưng rất được ưa chuộng.}{``Shortcuts are not evil if they come after understanding. But when tricks replace thinking, students get used to living by shortcuts. Over time they stop wanting to understand and only want to get through quickly. That is a poor form of achievement, yet a very popular one.''}
\end{truyen}

\section{Một Ngày Học Quá Dày Làm Não Bội Thực}
\begin{truyen}{Một Ngày Học Quá Dày Làm Não Bội Thực}{An Overloaded Day Gives the Brain Indigestion}
\chuhoa{C}{ó em học từ sáng tới tối, ngồi lớp này xong lao sang lớp kia.}
\textit{(Some students study from morning to night, leaving one class only to rush into another.)}

Đến tối về nhà chỉ còn cái xác biết ngồi.
\textit{(By nightfall only a body remains sitting at the desk.)}

\teacherpair{Não người không phải cái phễu muốn đổ bao nhiêu cũng được. Học cần hấp thụ, cần dừng, cần tiêu hóa. Nhồi mãi mà không có khoảng thở thì rất dễ biến việc học thành một dạng ép ăn tri thức tới mức sợ học luôn.}{``The human brain is not a funnel into which you can pour endlessly. Learning requires absorption, pause, and digestion. If you keep stuffing without breathing room, study itself starts to feel like force-feeding knowledge until the student grows afraid of learning.''}
\end{truyen}

\section{Con Đi Học Thêm Hay Đi Gánh Kỳ Vọng}
\begin{truyen}{Con Đi Học Thêm Hay Đi Gánh Kỳ Vọng}{Is the Child Studying, or Carrying Expectations?}
\chuhoa{N}{hiều đứa trẻ bước vào lớp thêm với chiếc cặp nặng, nhưng còn nặng hơn là ánh mắt mong chờ của người lớn.}
\textit{(Many children walk into tutoring class with a heavy bag, but heavier still are the adults' expectations.)}

\adultpair{Ba mẹ không sai khi muốn con tiến bộ. Nhưng nếu mọi buổi học thêm đều mang mùi kiểm soát, so sánh và sợ hãi, thì đứa trẻ không chỉ học kiến thức. Nó đang học luôn cảm giác mình chỉ đáng giá khi thành tích đủ đẹp.}{Parents are not wrong to want improvement. But if every tutoring session carries the smell of control, comparison, and fear, then the child is not only learning school material. The child is also learning to feel valuable only when the scores look good.}
\end{truyen}

\section{Muốn Học Thêm Có Ích, Phải Cắt Rác Trước}
\begin{truyen}{Muốn Học Thêm Có Ích, Phải Cắt Rác Trước}{If Tutoring Should Help, Cut the Junk First}
\chuhoa{H}{ọc thêm có thể có ích.}
\textit{(Tutoring can be useful.)}

Nhưng chỉ khi nó đi sau một cuộc chỉnh đốn thật.
\textit{(But only when it comes after honest correction.)}

\teacherpair{Ngủ cho đủ hơn. Bớt lướt đi. Tập ngồi yên lâu hơn. Làm bài sai rồi sửa lại. Dám hỏi chỗ mình mù. Khi cái nền này được dựng lại, học thêm mới có đất bám. Còn nếu đời sống vẫn nát, thêm lớp cũng chỉ là chồng thêm tiếng ồn lên một bộ não đã mệt.}{``Sleep more adequately. Scroll less. Practice sitting still longer. Redo mistakes. Dare to ask where you are blind. Once this foundation is rebuilt, tutoring finally has ground to stand on. But if daily life remains chaotic, adding more classes merely stacks more noise onto an already exhausted brain.''}
\end{truyen}

\begin{baihoc}
Học thêm không phải tội.
\textit{(Tutoring is not a sin.)}

Nhưng nếu dùng nó để che cho sự phân tán, thiếu tự học và lịch sống méo mó, thì nó chỉ làm vấn đề mặc áo tử tế hơn.
\textit{(But if it is used to cover up distraction, lack of self-study, and a crooked daily rhythm, it only dresses the problem in more respectable clothes.)}
\end{baihoc}

\begin{tuhocnhanh}
1. Em đang đi học thêm để hiểu hơn, hay chỉ để bớt hoảng? \textit{(Are you attending tutoring to understand more, or merely to panic less?)}

2. Trong tuần của em, giờ tự học thật có bị giờ học thêm nuốt mất không? \textit{(In your week, has real self-study been swallowed by tutoring time?)}

3. Nếu cắt bớt điện thoại và ngủ đủ hơn, một phần vấn đề học của em có tự giảm không? \textit{(If you reduced phone use and slept better, would part of your learning problem shrink on its own?)}
\end{tuhocnhanh}

\begin{tongketchuong}
Nhét thêm lớp vào một đời sống đang rối không bảo đảm tạo ra tiến bộ.
\textit{(Stuffing more classes into a disordered life does not guarantee progress.)}

Nhiều khi nó chỉ làm đứa trẻ trông có vẻ cố gắng hơn, trong khi bên trong vẫn đang rách dần.
\textit{(Sometimes it only makes the child look more hardworking while the inside continues to tear apart.)}
\end{tongketchuong}

\ngancach